\documentclass[11pt,a4paper]{article}

% --- Mathematics (must precede unicode-math) ---
\usepackage{amsmath}
\usepackage{mathtools}

% --- Fonts (lualatex) ---
\usepackage{fontspec}
\usepackage{unicode-math}
\setmathfont{KpMath-Regular.otf}
\setmathfont{KpMath-Bold.otf}[version=bold]

% --- Microtypography (after all font setup) ---
\usepackage{microtype}

% --- Colors & page layout ---
\usepackage[dvipsnames]{xcolor}
\usepackage[margin=25mm, top=30mm, bottom=30mm]{geometry}

% --- Headers & footers ---
\usepackage{fancyhdr}
\pagestyle{fancy}
\fancyhf{}
\fancyhead[L]{\small\itshape P--SV--SH Unitarity}
\fancyhead[R]{\small\itshape Nestor (2026)}
\fancyfoot[C]{\thepage}
\renewcommand{\headrulewidth}{0.4pt}
\renewcommand{\footrulewidth}{0pt}

% --- Section numbering ---
\setcounter{secnumdepth}{3}

% --- Tables ---
\usepackage{booktabs}
\usepackage{array}

% --- Captions ---
\usepackage[font=small, labelfont=bf, labelsep=period]{caption}

% --- Spacing ---
\usepackage{setspace}
\setstretch{1.15}
\setlength{\parskip}{0.4em}
\setlength{\emergencystretch}{3em}

% --- Bibliography ---
\usepackage[round,authoryear]{natbib}
\bibliographystyle{plainnat}

% --- Hyperlinks (last before macros) ---
\usepackage[colorlinks=true, linkcolor=blue!60!black,
            citecolor=blue!60!black, urlcolor=blue!60!black]{hyperref}

% --- Macros (house set) ---
\newcommand{\bG}{\mathbf{G}}
\newcommand{\bI}{\mathbf{I}}
\newcommand{\bx}{\mathbf{x}}
\newcommand{\dd}{\mathrm{d}}
\newcommand{\bT}{\mathbf{T}}

% --- New macros for this document ---
\newcommand{\bS}{\mathbf{S}}
\newcommand{\bR}{\mathbf{R}}
\newcommand{\ba}{\mathbf{a}}
\newcommand{\bb}{\mathbf{b}}
\newcommand{\RD}{\bR_{D}}
\newcommand{\RU}{\bR_{U}}
\newcommand{\TD}{\bT_{D}}
\newcommand{\TU}{\bT_{U}}
\newcommand{\etas}{\eta_{\beta}}
\newcommand{\kpar}{\mathbf{k}_{\parallel}}
\newcommand{\herm}{^{\dagger}}

\begin{document}

% --- Title block ---
\thispagestyle{plain}
\begin{center}
\vspace*{1em}
{\LARGE\bfseries P--SV--SH Unitarity:\\[0.3em]
Energy Flux, Open Channels and the Kennett Recursion}\\[1.8em]
{\large T.\,M.\ Nestor}\\[0.5em]
{\itshape Research School of Earth Sciences, Australian National University}\\[2em]
\end{center}

\noindent\textbf{Status.} The interface results of \S\ref{sec:numerics}
are validated by two independent implementations that agree to
$10^{-15}$: the Aki--Richards closed forms in \texttt{kennett\_layers.py}
(Python) and displacement--traction continuity in the thesis
energy-normalised eigenbasis, \texttt{InterfaceUnitarity.wl}
(Mathematica, 50 digits). The comparison is run by
\texttt{scripts/gate\_interface\_unitarity.py}. The periodic-plane
statement of \S\ref{sec:periodic} is the acceptance criterion of the
intra-plane energy-balance gate and is not re-derived here.

% =============================================================================
\section{The scattering matrix at fixed horizontal slowness}
\label{sec:smatrix}
% =============================================================================

A laterally homogeneous layer, or a stack of layers, preserves the
horizontal slowness $p$. At each $p$ and frequency $\omega$ it therefore
acts as a linear map from incoming to outgoing plane waves
\citep{Kennett1983,KennettKerry1979}. Write $\ba$ for the incoming
amplitudes (downgoing above, upgoing below) and $\bb$ for the outgoing
ones (upgoing above, downgoing below). Then
\begin{equation}
  \bb = \bS\,\ba,
  \qquad
  \bS =
  \begin{pmatrix}
    \RD & \TU \\
    \TD & \RU
  \end{pmatrix}.
  \label{eq:smatrix}
\end{equation}
In an isotropic medium the in-plane (P--SV) and anti-plane (SH) motions
decouple \citep[ch.~5]{AkiRichards2002}, so $\bS$ is block diagonal
with two independent pieces.
\begin{itemize}
\item \textbf{P--SV}: channels $\{\mathrm{P},\mathrm{SV}\}$ on each
  side. $\RD, \RU, \TD, \TU$ are $2\times2$ and $\bS$ is $4\times4$.
\item \textbf{SH}: one channel on each side. The blocks are scalars
  and $\bS$ is $2\times2$.
\end{itemize}
The unitarity statement
\begin{equation}
  \bS\herm \bS = \bI
  \label{eq:unitary}
\end{equation}
is therefore two statements, one for each block.

% =============================================================================
\section{Why flux normalisation is required}
\label{sec:flux}
% =============================================================================

Equation \eqref{eq:unitary} expresses energy conservation only if
$|a_k|^2$ \emph{is} the vertical energy flux carried by channel $k$.
For a plane wave of displacement amplitude $A$ the time-averaged
vertical flux is
\begin{equation}
  \langle F_z \rangle =
  \tfrac12 \omega^2 \, \rho \, v^2 \, \eta_v \, |A|^2,
  \qquad
  \eta_v = \sqrt{v^{-2} - p^2},
  \label{eq:flux}
\end{equation}
with $v = \alpha$ for P and $v = \beta$ for SV and SH
\citep[ch.~5]{AkiRichards2002}. For SH the prefactor is the vertical
impedance $Z = \rho\beta^2\etas$.

Displacement coefficients therefore obey an energy balance with
channel weights, for example
\begin{equation}
  1 = |R_{PP}|^2
    + \frac{\rho_1\beta_1^2\eta_{\beta1}}{\rho_1\alpha_1^2\eta_{\alpha1}}|R_{SP}|^2
    + \frac{\rho_2\alpha_2^2\eta_{\alpha2}}{\rho_1\alpha_1^2\eta_{\alpha1}}|T_{PP}|^2
    + \frac{\rho_2\beta_2^2\eta_{\beta2}}{\rho_1\alpha_1^2\eta_{\alpha1}}|T_{SP}|^2 .
  \label{eq:weighted}
\end{equation}
Rescaling each amplitude by the square root of its flux weight removes
the weights. The rescaled coefficients are the \emph{energy-normalised}
or \emph{modified} coefficients
\citep{Frasier1970,KennettKerry1979,Chapman2004}. In these variables,
\begin{equation}
  \sum_k |b_k|^2 = \sum_k |a_k|^2 \quad \text{for all } \ba
  \quad\Longleftrightarrow\quad
  \bS\herm\bS = \bI .
  \label{eq:equiv}
\end{equation}
The implementation uses $\sqrt{\rho\eta}$ factors for P--SV, with the
velocity factors carried by Kennett's amplitude convention, and
$\sqrt{Z}$ for SH.

\paragraph{SH example.} At a welded interface the modified SH
coefficients are
\begin{equation}
  R = \frac{Z_1 - Z_2}{Z_1 + Z_2}, \qquad
  T = \frac{2\sqrt{Z_1 Z_2}}{Z_1 + Z_2}, \qquad
  \bS_{\mathrm{SH}} =
  \begin{pmatrix} R & T \\ T & -R \end{pmatrix}.
  \label{eq:sh}
\end{equation}
For real $Z_1, Z_2$, $R^2 + T^2 = 1$ and $RT - TR = 0$, so
$\bS_{\mathrm{SH}}$ is real orthogonal and hence unitary.

% =============================================================================
\section{Content of the identity}
\label{sec:content}
% =============================================================================

Equation \eqref{eq:unitary} says that the columns of $\bS$ are
orthonormal. The two parts constrain different things.

\begin{itemize}
\item \textbf{Unit column norms: energy.} For an incident downgoing P
  wave, $|R_{PP}|^2 + |R_{SP}|^2 + |T_{PP}|^2 + |T_{SP}|^2 = 1$. This is
  \eqref{eq:weighted} with the weights absorbed.
\item \textbf{Orthogonal columns: phase.} For distinct inputs $j \neq
  l$, $\sum_k S_{kj}^{*} S_{kl} = 0$. Two different incident waves
  produce outgoing fields with no net interference flux. A model can
  have correct energies and still fail this condition if its phases are
  wrong.
\end{itemize}

% =============================================================================
\section{Conditions for unitarity}
\label{sec:conditions}
% =============================================================================

\begin{table}[h]
\centering
\caption{Conditions under which \eqref{eq:unitary} holds, and what
happens when each one fails.}
\label{tab:conditions}
\begin{tabular}{>{\raggedright}p{0.28\textwidth} >{\raggedright\arraybackslash}p{0.62\textwidth}}
\toprule
Condition & If violated \\
\midrule
Lossless media (real $\alpha$, $\beta$)
  & $\eta$ is complex, $\sqrt{\rho\eta}$ is no longer a flux
    normalisation, and \eqref{eq:unitary} carries no energy meaning.
    For a thick lossy layer propagation dominates and
    $\bS\herm\bS \preceq \bI$ approximately, the defect being absorbed
    energy. \\[0.4em]
Propagating channels only
  & Beyond a critical slowness, e.g.\ $p > 1/\alpha_2$, the channel is
    evanescent and carries no flux on its own. Remove it; the
    restriction of $\bS$ to open channels is unitary. \\[0.4em]
No open diffraction orders (periodic plane)
  & Energy leaves through channels with $\kpar + \mathbf{G} \neq
    \kpar$. The specular $\bS$ is sub-unitary even without loss. \\
\bottomrule
\end{tabular}
\end{table}

\paragraph{Evanescent channels.} An evanescent pair carries flux only
through the interference of its up- and downgoing parts. If such a
channel is kept, energy conservation is no longer $\bS\herm\bS = \bI$
but a statement in a non-diagonal flux metric
\citep{Newton1982}. Restricting to open channels avoids this.
Near the critical slowness $\eta \to 0$, so the normalisation factor
$\sqrt{\rho\eta}$ vanishes and the coefficients are poorly conditioned.

% =============================================================================
\section{Periodic planes: why specular unitarity is legitimate}
\label{sec:periodic}
% =============================================================================

A plane of cubic voxels with lattice pitch $a_L$ is laterally periodic,
not laterally invariant. Horizontal wavenumber is conserved only modulo
reciprocal-lattice vectors $\mathbf{G}$, and the unitary scattering
matrix lives in one Bloch sector. It includes \emph{every} open
channel $(\text{mode}, \pm, \mathbf{G})$ with real vertical
wavenumber
\begin{equation}
  k_z(\mathbf{G}) = \sqrt{\kappa^2 - |\kpar + \mathbf{G}|^2},
  \qquad \kappa \in \{\kappa_P, \kappa_S\}
  \label{eq:kz}
\end{equation}
\citep{Rayleigh1907,Petit1980}. When the plane is sub-wavelength, every
$\mathbf{G} \neq 0$ order is evanescent. Only the specular order is
open, and the $4\times4$ P--SV and $2\times2$ SH specular matrices are
unitary by themselves. The plane then acts as an effective laterally
homogeneous layer, which is the premise of an intra-plane Foldy--Lax
reduction to $R/T(p)$. If an order opens, specular unitarity is false,
and $\bI - \bS\herm\bS$ measures diffracted energy rather than error.
A numerical gate must refuse to run in that regime, not report a
defect.

% =============================================================================
\section{Unitarity versus reciprocity}
\label{sec:reciprocity}
% =============================================================================

These are independent properties.
\begin{itemize}
\item \textbf{Reciprocity} follows from the Betti--Rayleigh theorem.
  With Kennett's phase convention for the modified coefficients it
  makes $\bS$ symmetric: $\bS = \bS^{\mathsf T}$, i.e.\ $\RD =
  \RD^{\mathsf T}$, $\RU = \RU^{\mathsf T}$ and $\TU = \TD^{\mathsf T}$
  \citep{KennettKerry1979,VaratharajuluPao1976}. It is bilinear,
  contains no complex conjugate, and \emph{survives attenuation}.
  Its matrix form depends on the per-channel phase convention. In the
  thesis eigenbasis the P--SV block obeys instead
  \begin{equation}
    \bS^{\mathsf T} = \hat\Sigma\,\bS\,\hat\Sigma,\qquad
    \hat\Sigma = \mathrm{diag}(1,-1,1,-1),
    \label{eq:recip-thesis}
  \end{equation}
  i.e.\ $\RD^{\mathsf T} = \Sigma\RD\Sigma$, $\RU^{\mathsf T} =
  \Sigma\RU\Sigma$ and $\TU = \Sigma\TD^{\mathsf T}\Sigma$ with $\Sigma
  = \mathrm{diag}(1,-1)$, while the SH block is symmetric. Both hold to
  50 digits at an interface, with and without attenuation, and plain
  $\bS = \bS^{\mathsf T}$ fails there by $0.31$.
\item \textbf{Unitarity} is sesquilinear and holds \emph{only in the
  absence of loss}.
\end{itemize}
A lossless reciprocal system satisfies both. In the Kennett convention
$\bS^{-1} = \bS\herm = \bS^{*}$, which is the time-reversal statement that the conjugate
outgoing field retraces the incident one \citep{Waterman1976}. Because
the two properties constrain different combinations of entries, a
numerical check should test both.

% =============================================================================
\section{Unitarity of the Kennett recursion}
\label{sec:recursion}
% =============================================================================

The reflectivity recursion combines layer scattering matrices with the
Redheffer star product \citep{Redheffer1961,KennettKerry1979}. For two
elements $\bS^{(1)}$ above $\bS^{(2)}$,
\begin{equation}
  \RD^{(12)} = \RD^{(1)} + \TU^{(1)}\RD^{(2)}
    \bigl(\bI - \RU^{(1)}\RD^{(2)}\bigr)^{-1}\TD^{(1)},
  \label{eq:star}
\end{equation}
together with the analogous expressions for $\TD^{(12)}$, $\RU^{(12)}$
and $\TU^{(12)}$. The composite is the scattering matrix of the cascade.
Since each element conserves flux, the cascade does too, so the star
product of unitary matrices is unitary \citep{Pozar2012}. With modified
coefficients a lossless stack therefore remains unitary at every step,
and $\lVert \bS\herm\bS - \bI \rVert$ directly measures accumulated
numerical error. Displacement coefficients would require the flux
weights of \eqref{eq:weighted} to be carried through every step.

% =============================================================================
\section{Numerical check}
\label{sec:numerics}
% =============================================================================

The interface coefficients were evaluated for
$(\alpha, \beta, \rho)_1 = (3.0, 1.7, 2.3)$ and
$(\alpha, \beta, \rho)_2 = (4.2, 2.4, 2.7)$, in km/s and g/cm$^3$.
Attenuation uses the complex slowness
\begin{equation}
  s = \frac{1}{v}\,\frac{(2Q)^2 + \mathrm{i}\,2Q}{1 + (2Q)^2},
  \qquad \operatorname{Im} s > 0,
  \label{eq:qslow}
\end{equation}
with the branch $\operatorname{Re}\eta > 0$, $\operatorname{Im}\eta >
0$. The column $\lVert\cdot\rVert_{\max}$ reports the largest entry of
the residual.

\paragraph{Two independent routes.}
\begin{itemize}
\item \textbf{Python}, \texttt{kennett\_layers.py}: the Aki--Richards
  closed forms with $\sqrt{\rho\eta}$ and $\sqrt{Z}$ normalisation,
  in double precision.
\item \textbf{Mathematica}, \texttt{InterfaceUnitarity.wl}: continuity
  of the displacement--traction vector across $z = 0$ in the thesis
  energy-normalised eigenbasis $D_z$ (Nestor 1996, \S3.1), at 50 digits.
  With $\bb = D_1(\mathbf{d}_1;\mathbf{u}_1) = D_2(\mathbf{d}_2;\mathbf{u}_2)$
  and $Q = D_1^{-1}D_2$ in $3\times3$ blocks,
  \begin{equation}
    \TD = Q_{11}^{-1},\quad \RD = Q_{21}\TD,\quad \RU = -\TD Q_{12},\quad
    \TU = Q_{22} - Q_{21}\TD Q_{12}.
    \label{eq:qblocks}
  \end{equation}
\end{itemize}
The two share no code and no algebra. Both are flux-normalised but use
different per-channel phase conventions, so they are compared through
quantities invariant under unit-modulus rephasing: the magnitudes
$|S_{ij}|$ and the eigenvalues of $\bS\herm\bS$. As a negative control
the Python $\bS$ is converted back to displacement-basis coefficients,
$W^{-1/2}\bS W^{1/2}$ with $W = \rho v^2\eta$ from \eqref{eq:flux}. The
comparison must then fail, and it does, by $0.55$ to $1.4$ in every
case.

\begin{table}[h]
\centering
\caption{Unitarity residuals of the interface $\bS$ from the
Mathematica route, and its agreement with the Python route. The
$p = 0.30$ rows exceed $1/\alpha_2 \approx 0.238$, so P in medium 2 is
evanescent. ``0'' means zero to 50-digit working precision.}
\label{tab:numerics}
\small
\begin{tabular}{l l r}
\toprule
Case & Quantity & Value \\
\midrule
Lossless, $p \in \{0, 0.1, 0.2\}$ & $\lVert\bS\herm\bS - \bI\rVert_{\max}$, P--SV and SH & $0$ \\
Lossless, $p = 0.30$, all 4 channels & $\lVert\bS\herm\bS - \bI\rVert_{\max}$, P--SV & $1.56$ \\
Lossless, $p = 0.30$, 3 open channels & $\lVert\bS_r\herm\bS_r - \bI\rVert_{\max}$ & $0$ \\
Lossless, $p = 0.30$, SH & $\lVert\bS\herm\bS - \bI\rVert_{\max}$ & $0$ \\
$Q_P = 50$, $Q_S = 30$, $p = 0.1$ & eig$(\bS\herm\bS)$, P--SV & $0.9892$ to $1.0109$ \\
 & eig$(\bS\herm\bS)$, SH & $0.99949$, $1.00051$ \\
$Q_P = 20$, $Q_S = 10$, $p = 0.1$ & eig$(\bS\herm\bS)$, P--SV & $0.9666$ to $1.0346$ \\
 & eig$(\bS\herm\bS)$, SH & $0.99848$, $1.00152$ \\
\midrule
All six cases & Python vs Mathematica, $|S_{ij}|$, eig$(\bS\herm\bS)$ & $\le 1.0\times10^{-15}$ \\
All six cases & Displacement-basis control & $\ge 0.55$ \\
Lossy, Kennett basis & $\lVert\bS - \bS^{\mathsf T}\rVert_{\max}$ & $0$ \\
\bottomrule
\end{tabular}
\end{table}

The lossy rows show the point of \S\ref{sec:reciprocity}. Eigenvalues
of $\bS\herm\bS$ fall on \emph{both} sides of unity, and the spread
grows as $Q$ falls, so the lossy interface matrix is not a contraction
and \eqref{eq:unitary} has no energy interpretation there. Reciprocity
still holds exactly.

\bibliography{psvsh_unitarity}

\end{document}
